\documentclass[12pt]{article}
\usepackage[a4paper, total={6in, 8in}]{geometry}

\begin{document}


\textbf{PREFACE}

\hfill

Dynamical systems are an important topic in engineering. Applications are
prevalent within mechanical, electrical, and biomedical engineering and can
be found within robotic, automotive, aerospace, and human systems,
among others. The purpose of this book is to present dynamical systems
topics in a way that is relevant for practicing engineers. It is intended for
engineers who need to understand both the background theory and how to
apply it. As such, there is an attempt to create a bridge between the theory
and the application. Every abstract concept is discussed in depth, described
in a readable and down-to-earth manner, and illustrated using practical
examples. The intended audience is engineers who are working in industry,
graduate students who are taking courses or doing research related to
dynamical systems, and undergraduate students who are taking courses in
control systems. This is not a textbook, and there are no end-of-chapter
problems. Rather, it should be considered an application guide for those in
the trenches of working with and learning about dynamical systems.

It is assumed that readers have a solid mathematical foundation in calculus,
differential equations, and matrix theory. In presenting the material,
the emphasis is on applying the theory, so there are relatively few theorems
and no proofs. However, there is a lot of mathematics. Much mathematical
detail is given that is missing from other texts on these topics. The reason
for this level of detail is to help readers understand the complete application
in real-world systems. The focus is on depth and not breadth. In covering
the selected topics at this level of detail, unfortunately, the number of topics
had to be limited. As such, this is not a complete and comprehensive
presentation of all topics in dynamical systems. However, this book attempts
to cover many relevant topics that an engineer in the field would
encounter and provide a foundational understanding for further study.

It is also assumed that the reader has some understanding of MATLAB
and Simulink, although expertise is not required. Many of the concepts are
demonstrated using real-world examples in MATLAB or Simulink. For the
earlier chapters, the MATLAB code is explained line by line to show how
various concepts are implemented. These explanations are gradually
decreased throughout the book.

The book transitions from topics commonly found at the undergraduate
level in engineering, to those covered in graduate courses, to those that engineers
may never see in a course. As such, the coverage changes in its approach and
assumptions about what the reader knows. The layout of the topics is as follows.

Chapter 1 introduces dynamical systems, provides motivation for why it’s
important to study them, and discusses different types of systems. This material
should be familiar from undergraduate engineering courses in linear systems and
control theory. There is high-level discussion of this material, but the mathematics
starts early with definitions of the different classes of systems.

Chapter 2 discusses modeling and covers differential and difference
equations, transfer functions, state-space models, eigenvalues, eigenvectors,
and singular value decomposition. Although many of these topics are familiar
from courses in differential equations, control systems, and linear algebra, the
emphasis is on putting them in the context of dynamical systems. There is
also an emphasis on working through examples in MATLAB and giving
details of the implementation, which may not be covered in those courses.

Chapter 3 focuses on solutions of dynamical equations, equilibrium
points, and stability. These concepts are often encountered in introductory
graduate-level courses in dynamical systems and control theory. Again, the
emphasis is not on deriving the results but applying them. As such, several
of the relevant theorems are presented and applied.

Chapter 4 discusses nonlinear systems and some rich behavior that is
only found in them such as limit cycles, bifurcations, chaos, and linearization.
These are topics typically found in graduate-level engineering
courses. There is a minimal amount of theoretical coverage, and the behaviors
are described through examples.


Finally, Chapter 5 introduces Hamiltonian systems, which typically fall
in the realm of physicists. However, undamped vibrational systems and
their equivalent are an important class of Hamiltonian systems, and there is
much rich theory in this area. As with Chapter 4, there is minimal theoretical
coverage, and the focus is placed more on the introduction of the
concepts through examples.

Many people contributed to the creation of the book. Acknowledgement
goes out to colleagues and students at the University of Hartford,
particularly Lee Townsend and Iman Salehi for their supportive ideas and
engaging discussion; colleagues in the van Rooy Center for Complexity
and Conflict Analysis for the productive biweekly meetings; Harriet
Clayton and Glyn Jones at Elsevier for the support and feedback; and Joe
Romagnano for his editorial skills.

\begin{thebibliography}{9}
\bibitem{texbook}
PATRICIA MELLODGE (2016) \emph{A Practical
Approach to
Dynamical Systems
for Engineers}, Woodhead Publishing Series.

\bibitem{lamport94}
Leslie Lamport (1994) \emph{\LaTeX: a document preparation system}, Addison
Wesley, Massachusetts, 2nd ed.
\end{thebibliography}

\end{document}


